\begin{definition}[BNT operators and structure coefficients]%
    \label{def:mpdo_bnt_label_theorem_witness}
    \lean{MPOTensor.BNTLabelTheoremWitness}
    \lean{MPOTensor.HasBNTLabelTheoremWitness}
    \lean{MPOTensor.BNTLabelTheoremWitness.toTheoremData}
    \lean{MPOTensor.BNTLabelTheoremWitness.coeffs}
    \lean{MPOTensor.BNTLabelTheoremWitness.operators}
    \lean{MPOTensor.BNTLabelTheoremWitness.traceScalars}
    \lean{MPOTensor.BNTLabelTheoremWitness.positiveChi}
    \lean{MPOTensor.BNTLabelTheoremWitness.blockedComparison}
    \lean{MPOTensor.BNTLabelTheoremWitness.same_length_product_form}
    \lean{MPOTensor.BNTLabelTheoremWitness.idempotent_coefficient_form}
    \lean{MPOTensor.BNTLabelTheoremWitness.same_length_product_form_ofChi}
    \lean{MPOTensor.BNTLabelTheoremWitness.idempotent_coefficient_form_ofChi}
    \lean{MPOTensor.BNTLabelTheoremWitness.same_length_product_eq_sum_ofChi}
    \lean{MPOTensor.BNTLabelTheoremWitness.idempotent_eq_sum_ofChi}
    \lean{MPOTensor.BNTLabelTheoremWitness.positive_chi_pos_entries}
    \lean{MPOTensor.BNTLabelTheoremWitness.positive_chi_trace_power}
    \lean{MPOTensor.BNTLabelTheoremWitness.labelAssignment}
    \lean{MPOTensor.BNTLabelTheoremWitness.sourceLabel}
    \lean{MPOTensor.BNTLabelTheoremWitness.targetLabel}
    \lean{MPOTensor.BNTLabelTheoremWitness.positiveBlockedChi}
    \lean{MPOTensor.BNTLabelTheoremWitness.ofChi}
    \lean{MPOTensor.BNTLabelTheoremWitness.blocked_coeff_eq}
    \lean{MPOTensor.BNTLabelTheoremWitness.blockedComparison_ofChi}
    \lean{MPOTensor.BNTLabelTheoremWitness.blocked_coeff_eq_ofChi}
    \leanok
    \uses{def:mpdo_bnt_label_theorem_data}
    Fix the algebra-structure data. Consider a finite set of labels $\alpha$,
    complex vector spaces equipped with a product at each length $L$, and
    operators $O_L(M_\alpha)$ in those spaces, together with scalars
    $m_\alpha$, structure coefficients $c^{(L)}_{\alpha,\beta,\gamma}$,
    and length-independent diagonal matrices
    $\chi_{\alpha,\beta,\gamma}$ with positive diagonal entries. Require
    the product, idempotent, trace-power, and blocked-basis comparison
    conditions of Definition~\ref{def:mpdo_bnt_label_theorem_data}.

    In a vertical BNT decomposition, $O_L(M_\alpha)$ denotes the operator
    obtained by contracting $L$ copies of $M_\alpha$, and
    $m_\alpha=\tr(\mu_\alpha)$ for the positive diagonal multiplicity
    matrix $\mu_\alpha$. The abstract conditions here do not themselves
    construct such a tensor decomposition.
    % The converse direction of CPSV16 Theorem 4.14 compares the algebra-side
    % and blocked physical sectors through marked/common-target tensors, proves
    % X_\gamma^\dagger X_\gamma=\omega_\gamma I with \omega_\gamma>0, normalizes
    % X_\gamma to U_\gamma=\omega_\gamma^{-1/2}X_\gamma, and constructs the
    % physical tpCP maps T=\widetilde R_2\widetilde T R_1 and
    % S=\widetilde R_1\widetilde S R_2. The mixed-prefix comparison supplies
    % the common target under the standing canonical-form and MPDO hypotheses.
    Existence means that these objects can be chosen simultaneously.
    The blocked-basis comparisons are additional conditions relating them
    to the fixed algebra-structure data, not clauses of the BNT algebra
    characterization.

    One may instead start with the positive diagonal matrices and define
    $c^{(L)}_{\alpha,\beta,\gamma}
        =\tr(\chi_{\alpha,\beta,\gamma}^{\,L})$,
    provided the product, idempotent, and blocked-basis comparison
    conditions hold. For any choice satisfying the conditions above,
    replacing the coefficients by these traces at positive lengths
    preserves all three conditions.
\end{definition}

\begin{theorem}[Product and idempotent equations for BNT operators]%
    \label{thm:mpdo_has_bnt_label_theorem_witness_source_equations}
    \lean{MPOTensor.HasBNTLabelTheoremWitness.exists_source_chi_trace_equations}
    \lean{MPOTensor.BNTLabelTheoremWitness.coeff_eq_trace_pow}
    \lean{MPOTensor.BNTLabelTheoremWitness.coeff_eq_ofChi_coeff}
    \lean{MPOTensor.BNTLabelTheoremWitness.chi_entry_pos}
    \lean{MPOTensor.BNTLabelTheoremWitness.same_length_product_eq_sum}
    \lean{MPOTensor.BNTLabelTheoremWitness.idempotent_eq_sum}
    \lean{MPOTensor.BNTLabelTheoremWitness.same_length_product_eq_sum_chi_trace_pow}
    \lean{MPOTensor.BNTLabelTheoremWitness.idempotent_eq_sum_chi_trace}
    \leanok
    \uses{def:mpdo_bnt_label_theorem_witness}
    Suppose the objects in
    Definition~\ref{def:mpdo_bnt_label_theorem_witness} exist.
    For any such choice, the product law holds for every positive length
    $L$ and all labels $\alpha,\beta$, and the scalar identity holds for
    every label $\gamma$:
    \begin{align}
        O_L(M_\alpha)O_L(M_\beta)
         & = \sum_\gamma c^{(L)}_{\alpha,\beta,\gamma} O_L(M_\gamma),
        \label{eq:mpdo_bnt_same_len_product}                                    \\
        m_\gamma
         & = \sum_{\alpha,\beta} c^{(1)}_{\alpha,\beta,\gamma}m_\alpha m_\beta.
        \label{eq:mpdo_bnt_idempotent_law}
    \end{align}
    Thus the vector $m=(m_\alpha)_\alpha$ is an idempotent for the
    multiplication induced by $c^{(1)}$.
    The diagonal matrices $\chi_{\alpha,\beta,\gamma}$ have positive
    diagonal entries and satisfy
    $c^{(L)}_{\alpha,\beta,\gamma}
        = \tr(\chi_{\alpha,\beta,\gamma}^{\,L})$
    for every $L>0$.
    Substituting this identity into~(\ref{eq:mpdo_bnt_same_len_product})
    and~(\ref{eq:mpdo_bnt_idempotent_law}) gives
    \begin{align}
        O_L(M_\alpha)O_L(M_\beta)
         & = \sum_\gamma \tr(\chi_{\alpha,\beta,\gamma}^{\,L}) O_L(M_\gamma),
        \label{eq:mpdo_bnt_same_len_product_chi}                                   \\
        m_\gamma
         & = \sum_{\alpha,\beta} \tr(\chi_{\alpha,\beta,\gamma}) m_\alpha m_\beta.
        \label{eq:mpdo_bnt_idempotent_law_chi}
    \end{align}
    At positive lengths, the coefficients therefore agree with those
    defined directly by the traces of powers of $\chi$.

    Let $\sigma_n$ and $\tau_n$ be the source and target label maps in
    the additional blocked-basis comparison at a positive length $n$.
    For source indices $i,j$ at length $n$ and a target index $k$ at
    length $2n$, the blocked-basis coefficients satisfy
    \begin{align}
        c^{(n)}_{i,j,k}
         & = c^{(n)}_{\sigma_n(i),\sigma_n(j),\tau_n(k)}
        = \tr\!\left(\chi_{\sigma_n(i),\sigma_n(j),\tau_n(k)}^{\,n}\right).
        \notag
    \end{align}
\end{theorem}
\begin{proof}\leanok
    Choose the assumed operators, scalars, coefficients, and diagonal
    matrices. Their product and idempotent identities give
    (\ref{eq:mpdo_bnt_same_len_product})
    and~(\ref{eq:mpdo_bnt_idempotent_law}); substitute the trace-power
    identity to obtain
    (\ref{eq:mpdo_bnt_same_len_product_chi})
    and~(\ref{eq:mpdo_bnt_idempotent_law_chi}).
    The blocked-basis comparison gives the final coefficient formula.
\end{proof}

\begin{theorem}[Positive diagonal matrices for blocked coefficients]%
    \label{thm:mpdo_bnt_label_theorem_witness_to_positive_blocked_chi_trace_power_form}
    \lean{MPOTensor.BNTLabelTheoremWitness.toPositiveBlockedStructureChiTracePowerForm}
    \lean{MPOTensor.BNTLabelTheoremWitness.positiveBlockedChi_toDiagonal_of_pos}
    \lean{MPOTensor.BNTLabelTheoremWitness.positiveBlockedChi_tracePowerCoeff_of_pos}
    \lean{MPOTensor.BNTLabelTheoremWitness.positiveBlockedChi_dim_of_pos}
    \lean{MPOTensor.BNTLabelTheoremWitness.positiveBlockedChi_trace_matrix_pow_of_pos}
    \lean{MPOTensor.BNTLabelTheoremWitness.positiveBlockedChi_tracePower}
    \lean{MPOTensor.BNTLabelTheoremWitness.positiveBlockedChi_posEntries}
    \lean{MPOTensor.BNTLabelTheoremWitness.positiveBlockedChi_entry_pos}
    \lean{MPOTensor.HasBNTLabelTheoremWitness.positive_blocked_chi_witness}
    \lean{MPOTensor.HasBNTLabelTheoremWitness.exists_blocked_chi_trace_power_form}
    \lean{MPOTensor.HasBNTLabelTheoremWitness.exists_blocked_coeff_eq_trace_pow}
    \leanok
    \uses{def:mpdo_bnt_label_theorem_witness, thm:mpdo_bnt_label_theorem_data_to_positive_blocked_chi_trace_power_form}
    Suppose the objects in
    Definition~\ref{def:mpdo_bnt_label_theorem_witness} exist for the
    chosen algebra-structure data. Then there are diagonal matrices
    $\chi_{n,i,j,k}$ with positive diagonal entries such that, for every
    positive blocked length $n$, source indices $i,j$ at length $n$,
    and target index $k$ at length $2n$,
    \begin{align}
        \chi_{n,i,j,k}
         & = \chi_{\sigma_n(i),\sigma_n(j),\tau_n(k)},
        \notag                                         \\
        c^{(n)}_{i,j,k}
         & = \tr(\chi_{n,i,j,k}^{\,n}).
        \notag
    \end{align}
    These matrices are obtained by selecting the BNT matrices
    $\chi_{\alpha,\beta,\gamma}$ along the source and target comparison
    maps $\sigma_n,\tau_n$. Their dimensions, diagonal entries, and
    traces of powers agree with those of the selected matrices.
    The same construction proves existence when the BNT operators,
    scalars, and matrices are assumed to exist without a fixed choice.
\end{theorem}
\begin{proof}\leanok
    Apply the blocked-basis trace-power consequence to the chosen
    operators, scalars, coefficients, and diagonal matrices.
\end{proof}

\begin{corollary}[Blocked coefficients as traces of powers]%
    \label{cor:mpdo_bnt_label_theorem_witness_blocked_coeff_eq_positive_blocked_chi_trace_pow}
    \lean{MPOTensor.BNTLabelTheoremWitness.blocked_coeff_eq_trace_pow}
    \lean{MPOTensor.BNTLabelTheoremWitness.blocked_coeff_eq_positiveBlockedChi_trace_pow}
    \leanok
    \uses{def:mpdo_bnt_label_theorem_witness, thm:mpdo_bnt_label_theorem_witness_to_positive_blocked_chi_trace_power_form, lem:mpdo_blocked_structure_chi_eq_trace_matrix_pow}
    Assume the conditions of
    Definition~\ref{def:mpdo_bnt_label_theorem_witness}.
    For every positive blocked length $n$, source indices $i,j$ at
    length $n$, and target index $k$ at length $2n$, the blocked-basis
    coefficients are traces of powers of the BNT matrices selected by
    the source and target comparison maps:
    \begin{align}
        c^{(n)}_{i,j,k}
         & = \tr\!\left(\chi_{\sigma_n(i),\sigma_n(j),\tau_n(k)}^{\,n}\right)
        = \tr(\chi_{n,i,j,k}^{\,n}).
        \notag
    \end{align}
\end{corollary}
\begin{proof}\leanok
    Apply the trace-power identity for the positive diagonal matrices
    $\chi_{n,i,j,k}$ obtained from the assumed BNT matrices and
    blocked-basis comparison.
\end{proof}
